\chapter{Glue Ear}

\section*{Definition and Overview}
Glue ear, medically known as otitis media with effusion, is a common condition in which thick, sticky fluid accumulates in the middle ear behind the eardrum without signs of acute infection. It is most common in young children, who may not complain of symptoms but can have reduced hearing, which may affect speech, language, and learning. Glue ear often follows colds or ear infections and usually resolves on its own within a few months. When fluid persists for a long time or affects development, treatment options include careful monitoring, hearing aids in selected cases, or surgery to insert ventilation tubes (grommets).

\section*{Epidemiology}
Glue ear is very common in childhood. Around 1 in 5 children have glue ear at any given time, and most children will have at least one episode before starting school. It peaks between the ages of 1 and 6, with another smaller peak in early adolescence. Boys are slightly more affected than girls, and risk is higher in children who attend childcare, who are exposed to cigarette smoke, who have frequent colds, or who have a family history. Most cases resolve spontaneously, and fewer than 1 in 10 children need surgery.

\section*{Aetiology and Pathophysiology}
Glue ear develops when the middle ear does not drain properly and fluid collects behind the eardrum.

\begin{itemize}
    \item \textbf{Eustachian tube dysfunction:} the tube connecting the middle ear to the back of the nose is short and poorly supported in young children, so it does not drain effectively
    \item \textbf{Triggering events:} colds, viral infections, allergies, and acute ear infections cause swelling of the tube and impair drainage
    \item \textbf{Fluid accumulation:} the middle ear fills with fluid, which may be thin or thick and "glue-like"
    \item \textbf{Negative pressure:} impaired ventilation of the middle ear creates suction that draws in more fluid
    \item \textbf{Conductive hearing loss:} the fluid dampens the vibration of the eardrum, reducing hearing, typically by 10--30 decibels
    \item \textbf{Persistence:} in some children the fluid persists for months, affecting speech and learning
\end{itemize}

\section*{Clinical Features}
Glue ear is often symptomless and detected at routine hearing checks or school screening.

\begin{itemize}
    \item Mild to moderate hearing loss, which may be noticed as the child not responding, turning up the television, or speaking loudly
    \item Children appearing inattentive or struggling at preschool and school
    \item Speech and language delay in younger children
    \item A feeling of fullness or popping in the ear, and sometimes mild discomfort
    \item Balance problems and clumsiness in some children
    \item Earache is not typical, unlike acute otitis media
    \item The eardrum may look dull or retracted with fluid visible behind it
\end{itemize}

\section*{Differential Diagnosis}
Other causes of hearing problems and ear symptoms should be considered.

\begin{itemize}
    \item Acute otitis media: a painful ear infection with fever and a red, bulging eardrum
    \item Chronic suppurative otitis media: persistent discharge through a perforated eardrum
    \item Sensorineural hearing loss from other causes
    \item Impacted earwax, which can also cause conductive hearing loss
    \item Foreign body in the ear
    \item Speech and language delay from developmental or other causes
    \item Auditory processing disorder, which can mimic the effects of glue ear
\end{itemize}

\section*{When to Seek Medical Care}
See a doctor if you suspect a child has hearing problems, poor attention, or speech delay, or if ear symptoms persist. Children with suspected glue ear should have a formal hearing assessment, and referral to an ear, nose, and throat (ENT) specialist may be arranged for persistent cases.

\textbf{Call 000 (or local emergency services) for:}
\begin{itemize}
    \item This condition is not an emergency; however, seek urgent review for severe ear pain, high fever, or discharge from the ear, which may indicate an acute infection
    \item Sudden hearing loss in an older child or adult needs prompt medical assessment
\end{itemize}

\section*{Diagnosis and Investigations}
Glue ear is diagnosed by examining the eardrum and confirmed with hearing tests.

\begin{itemize}
    \item \textbf{Otoscopy:} examination of the eardrum, which may look dull, retracted, or amber-coloured with fluid visible behind it
    \item \textbf{Tympanometry:} a quick test that measures the movement of the eardrum and detects fluid in the middle ear
    \item \textbf{Hearing assessment:} formal audiometry, adapted to the child's age, measures the degree of hearing loss
    \item \textbf{Observation over time:} repeated checks determine whether the fluid is resolving
    \item \textbf{Referral:} ENT review is arranged for persistent or severe cases, or when there are concerns about development
\end{itemize}

\section*{Management}
Most children with glue ear do not need treatment beyond monitoring, and most cases resolve spontaneously.

\begin{itemize}
    \item \textbf{Watchful waiting:} observation for three months is standard, as most cases resolve on their own
    \item \textbf{Review of hearing:} repeat hearing tests determine whether hearing loss is significant or persistent
    \item \textbf{Communication strategies:} speaking clearly, facing the child, and reducing background noise help children with mild hearing loss
    \item \textbf{Managing contributing factors:} treating allergies, avoiding second-hand smoke, and reducing exposure to infections
    \item \textbf{Grommets (ventilation tubes):} if fluid persists for more than three months with significant hearing loss, small tubes are inserted through the eardrum to ventilate the middle ear
    \item \textbf{Adenoidectomy:} removal of enlarged adenoids may be recommended in some children, especially older ones
    \item \textbf{Hearing aids:} an option in children with persistent glue ear who are not suitable for or awaiting surgery
\end{itemize}

\section*{Complications}
\begin{itemize}
    \item Persistent conductive hearing loss and its effects on speech, language, and learning
    \item Behavioural and attention problems related to hearing difficulty
    \item Recurrent acute ear infections in some children
    \item Tympanosclerosis: scarring of the eardrum, which may occur after grommets but rarely affects hearing
    \item Grommet-related problems, including discharge, blockage, or the need for reinsertion
    \item Long-term educational effects if persistent hearing loss is not addressed
\end{itemize}

\section*{Prognosis and Outcomes}
The outlook for glue ear is excellent. About half of all cases resolve within three months, and most within a year, as the Eustachian tube matures. Even in children who need grommets, hearing returns to normal immediately after insertion, and most children require only one set. Speech and learning usually catch up once hearing is restored. Children with persistent glue ear and unrecognised hearing loss are at risk of educational delay, which is why routine hearing checks and early referral matter.

\section*{Prevention and Self-Care}
\begin{itemize}
    \item Avoid exposing children to second-hand cigarette smoke
    \item Practise good hand hygiene to reduce the frequency of colds
    \item Ensure immunisations are up to date, including the pneumococcal vaccine
    \item Breastfeed if possible, as breastfeeding is associated with fewer ear infections
    \item Watch for signs of hearing difficulty in young children and seek assessment early
    \item Speak clearly and face the child when talking, and reduce background noise
    \item Attend follow-up hearing tests as arranged
    \item Discuss any concerns about speech, language, or school progress with your doctor
\end{itemize}

\section*{Sources and Further Reading}
The following reputable sources provide further information for healthcare students and patients seeking reliable, current advice about glue ear.

\begin{itemize}
    \item Healthdirect Australia. \textit{Glue ear.} https://www.healthdirect.gov.au/
    \item Australian Society of Otolaryngology Head and Neck Surgery. \textit{Patient information on glue ear.} https://www.asohns.org.au/
    \item Raising Children Network (Australian Government). \textit{Glue ear and ear infections.} https://raisingchildren.net.au/
    \item NHS. \textit{Glue ear.} https://www.nhs.uk/
    \item Mayo Clinic. \textit{Ear infection and fluid in the ear.} https://www.mayoclinic.org/
    \item MSD Manual. \textit{Otitis media with effusion.} https://www.msdmanuals.com/
\end{itemize}
